\documentclass[10pt,a4paper]{amsart}
\usepackage[margin=19mm]{geometry}
\usepackage{amsmath,amssymb,mathtools}
\usepackage[scr=rsfs,cal=euler]{mathalfa}
\usepackage{xfrac}
\usepackage[unicode,hidelinks,bookmarks=false]{hyperref}
\newcommand{\ie}{i.e.\ }
\newcommand{\cf}{cf.\ }
\newcommand{\resp}{respectively\ }
\newcommand{\Subsection}{\subsection}
\providecommand{\nobreakdash}{\nobreak\hbox{-}}
\setlength{\emergencystretch}{2em}
\multlinegap0pt
\usepackage{amssymb}
\usepackage{float}
\usepackage{mathtools}
\usepackage{tikz}
\newtheorem{lemma}{Lemma}
\title{On Ulam's Segment Motion Problem}
\author{Dragomir Grozev, Nikolai Nikolov}
\date{Source: arXiv:2606.23268v1 (CC BY 4.0)}
\begin{document}
\maketitle

\noindent\emph{Source and licence.} On Ulam's Segment Motion Problem. Version arXiv:2606.23268v1 (CC BY 4.0), distributed by arXiv under the Creative Commons Attribution 4.0 International licence, http://creativecommons.org/licenses/by/4.0/, as recorded in the arXiv licence metadata for this exact version and shown on the abstract page https://arxiv.org/abs/2606.23268v1. Full licence notice: \texttt{licenses/arxiv-2606.23268-CC-BY-4.0.txt}.\par\smallskip

\noindent\textit{Editorial scope.} Counted unit: the article's lemma lem:min\_path (source0.tex lines 597--626). The article's complete original proof of it is delivered with it (source0.tex lines 628--693, one \texttt{proof} environment of 66 lines). No mathematics is written, repaired or re-derived: the statement and the proof are the article's own lines, in the article's own notation, and no retained line is substituted or re-flowed.  No further source-local context is needed: every cross-reference of every retained line resolves inside the two delivered spans.  Nothing was omitted from any retained span, so the package carries no omission record.  No retained line contains a citation, so the wrapper carries no bibliography and no bibliography entry.  No retained line contains a figure environment, an image-inclusion command or drawing code, so no image is delivered and the package declares no asset dependency.  Editorial formatting only, in addition: an AMS \texttt{amsart} preamble that restates the article's own declarations and macros used by the retained lines, namely the environments \texttt{lemma} on separate counters, together with the standard packages those lines need.  The retained environments print in this document as Lemma 1 in order of appearance, whereas the article prints the same items with the numbering of its own full document.  The complete native source of the article as distributed in the arXiv e-print for this exact version is bundled unchanged as \texttt{sources/203418/source0.tex}.
\begin{lemma}
\label{lem:min_path}
Let \(u_0,u_1\in\mathbb R^{d-1}\), $d\ge 3$, satisfy
\[
|u_0|=|u_1|=1,
\]
and let \(\theta\), \(0\le \theta\le \pi\), be the angle between \(u_0\) and \(u_1\).
Let \(O\) be the origin, and $0\le r\le 1$.

Among all rectifiable curves \(\gamma\subset\mathbb R^{d-1}\) joining \(u_0\) to \(u_1\)
and satisfying
\[
|x-O|\ge r
\qquad\text{for every }x\in\gamma,
\]
the minimum possible length is
\[
\ell_{\min}
=
\begin{cases}
2\sin\dfrac{\theta}{2},
&
\theta\le 2\arccos r,
\\[8pt]
2\sqrt{1-r^2}+r\bigl(\theta-2\arccos r\bigr),
&
\theta\ge 2\arccos r.
\end{cases}
\]
\end{lemma}

\begin{proof}
First, we show that the minimum is attained by a curve contained in the plane spanned by \(u_0\) and \(u_1\). 
Let \(\tilde u(\varphi)\) be the unit vector in the plane spanned by \(u_0\) and \(u_1\) such that
\[
\angle(u_0,\tilde u(\varphi))=\varphi,
\qquad
\varphi\in[0,\theta].
\]
Let \(x=x(t)\in\mathbb R^{d-1}\) be a curve avoiding the ball of radius \(r\) centered at the origin. We write
\[
x(t)=\rho(t)u(t),
\]
where \(\rho(t)\ge r\) and \(|u(t)|=1\). Reparametrizing the curve if necessary, we may assume that \(u(t)\) is parametrized by spherical arc length, that is,
\[
s=\int_0^t |\dot u(\tau)|\,d\tau,
\qquad t\in[0,1].
\]
Thus, $x(s)=\rho(s)u(s)$, $s\in[0,\ell]$, where $\rho(s)\ge r$, $u(s)\in \mathbb{R}^{d-1}$, $|u(s)|=1$ and $\displaystyle \ell:=\int_0^1 |\dot{u}|\,dt$. Consider the following planar curve
\[
\tilde{x}(s):= \rho(s)\tilde{u}\left(\frac{\theta}{\ell} s\right).
\]
It also avoids the ball $\{x:|x|<r\}$ and lies in the plane spanned by $u_0,u_1$.
We have
\[
d\tilde{x} =d \rho \tilde{u}+\rho \frac{\theta}{\ell}d\tilde{u},\quad dx =d \rho u+\rho du.
\]
Since $|u|=|\tilde{u}|=1$ and $u\perp du$, $\tilde{u}\perp d \tilde{u}$, $|u'(s)|=1$, $|\tilde{u}'(\varphi)|=1$, we get
\[
|d\tilde{x}|^2=d\rho^2|\tilde{u}|^2+\rho^2 \frac{\theta^2}{\ell^2}|d\tilde{u}|^2=d\rho^2+\rho^2\frac{\theta^2}{\ell^2}\,|ds|^2
\]
\[
|dx|^2=d\rho^2+\rho^2\,|ds|^2
\]
Since $\displaystyle \theta/\ell\le 1$, this means $|d\tilde{x}|\le |dx|$. Therefore
\[
\int_0^{\ell}|\tilde{x}'(s)|\,ds\le \int_0^{\ell}|x'(s)|\,ds 
\]

Thus, we have replaced the curve $x(t)$ by a planar one $\tilde{x}(t)$ of no greater length. Therefore, it suffices to consider the planar problem. The endpoints lie on the unit circle and have angular separation \(\theta\), while the curve is constrained to remain outside the disk of radius \(r\) centered at the origin.

If the chord joining \(u_0\) and \(u_1\) does not intersect the open disk \(|x|<r\),
then the chord is admissible and is the shortest path. The distance from the origin
to this chord is \(\cos(\theta/2)\). Hence this case occurs precisely when $\cos(\theta/2)\ge r$, or equivalently,
\[
\theta\le 2\arccos r.
\]
The minimum length is then
\[
|u_0-u_1|=2\sin\frac{\theta}{2}.
\]
Assume now that $\theta>2\arccos r$.
Then the chord intersects the forbidden disk. The shortest admissible curve consists
of a tangent segment from \(u_0\) to the circle \(|x|=r\), followed by the shorter
arc of this circle, followed by a tangent segment from the circle \(|x|=r\) to \(u_1\).

Each tangent segment has length $\sqrt{1-r^2}$. If \(\alpha\) is the angle between \(u_0\) and the radius to the tangent point, then $\cos\alpha=r$, and hence $\alpha=\arccos r$. Therefore the central angle of the circular arc is
\[
\theta-2\alpha=\theta-2\arccos r. 
\]
The length of the circular arc is $r\bigl(\theta-2\arccos r\bigr)$. Thus in this case the minimum length is
\[
2\sqrt{1-r^2}+r\bigl(\theta-2\arccos r\bigr).
\]

The two expressions coincide when \(\theta=2\arccos r\), so the proof is complete.
\end{proof}

\end{document}
